\chapter{What the paper did}
\label{ch:paper}

Everything in this chapter belongs to Lu, Kaplan \& Buehler. My own work starts
in Chapter~\ref{ch:reproduction}.

\section{The three tasks}

Fine-tuning taught the model three text formats. The formats \emph{are} the
interface --- there is no other API.

\begin{figure}[htbp]
\centering
\begin{tikzpicture}[font=\small,
  lab/.style={font=\scriptsize\bfseries,text=papercol},
  tape/.style={draw=greyc!70,fill=codebg,rounded corners=1.5pt,inner sep=5pt,
               font=\ttfamily\scriptsize},
]
\node[lab] (l1) at (0,2.3) {1. \texttt{Sequence} --- used in pretraining};
\node[tape,below=2pt of l1.south west,anchor=north west,text width=12.5cm]
  {Sequence<GGDSGAAAAAAAADGGRGGYGGLGRGGDSGAAAAAAAADGG\ldots>};

\node[lab] (l2) at (0,0.9) {2. \texttt{CalculateSilkContent} --- the \emph{forward} task};
\node[tape,below=2pt of l2.south west,anchor=north west,text width=12.5cm]
  {CalculateSilkContent<AAAGGAGQGGYGGQGAGQGA\ldots>\\
   \textcolor{ourscol}{[0.327,0.356,0.261,0.287,0.437,0.190,0.220,0.301]}};

\node[lab] (l3) at (0,-0.9) {3. \texttt{GenerateSilkContent} --- the \emph{inverse} task};
\node[tape,below=2pt of l3.south west,anchor=north west,text width=12.5cm]
  {GenerateSilkContent<0.327,0.356,0.261,0.287,0.437,0.190,0.220,0.301>\\
   \textcolor{ourscol}{[AAAGGAGQGGYGGQGAGQGAAAAAAGGAGQGGYGG\ldots]}};
\end{tikzpicture}
\caption[The three tasks]{The three fine-tuning formats. Black is the prompt
you supply; orange is what the model generates. The angle brackets hold the
input, the square brackets hold the output. Note that tasks 2 and 3 are the
same information in opposite directions --- that symmetry is what makes the
round trip of Section~\ref{sec:cycle} possible.}
\label{fig:tasks}
\end{figure}

\begin{keybox}{The eight-dimensional property vector}
Equation~(1) of the paper fixes the order, and nothing in the prompt labels the
slots --- the model learned them positionally:
\[
  [\,\underbrace{\text{toughness}}_{0},\;
    \underbrace{\text{SD}}_{1},\;
    \underbrace{E}_{2},\;
    \underbrace{\text{SD}}_{3},\;
    \underbrace{\text{strength}}_{4},\;
    \underbrace{\text{SD}}_{5},\;
    \underbrace{\text{strain}}_{6},\;
    \underbrace{\text{SD}}_{7}\,]
\]
All eight are \term{min--max normalised} to $[0,1]$: for each property, the
minimum and maximum across the dataset are found, and every value is mapped by
\[
  v_{\text{norm}} = \frac{v_{\text{raw}} - \min}{\max - \min}.
\]
Values are always written to three decimal places. Recovering the eight
$(\min,\max)$ pairs, which live in a supplementary table I could not initially
obtain, became one of the more satisfying parts of the project
(Section~\ref{sec:normalisation}).
\end{keybox}

\section{The cycle-consistency procedure}
\label{sec:cycle}

This is the procedure that produces every headline number.

\begin{figure}[htbp]
\centering
\begin{tikzpicture}[font=\small,
  n/.style={draw=greyc!70,fill=boxbg,rounded corners=2pt,inner sep=6pt,
            text width=2.9cm,align=center,minimum height=1.0cm},
  pn/.style={n,draw=papercol!60,fill=papercol!8},
  on/.style={n,draw=ourscol!60,fill=ourscol!8},
  ar/.style={-{Stealth[length=2.4mm]},thick,greyc!80},
]
\node[pn] (target) {\textbf{target}\\ 8 numbers you want};
\node[n,right=1.4cm of target] (gen) {\texttt{GenerateSilkContent}\\ \scriptsize sample $N$ times};
\node[n,below=1.0cm of gen] (filter) {\textbf{filter}\\ \scriptsize must parse as protein;\\ must be novel};
\node[n,left=1.4cm of filter] (fwd) {\texttt{CalculateSilkContent}\\ \scriptsize on each survivor};
\node[on,below=1.0cm of fwd] (score) {\textbf{score}\\ $R^2$(target, predicted)};
\node[on,right=1.4cm of score] (max) {\textbf{report}\\ the maximum};

\draw[ar] (target) -- (gen);
\draw[ar] (gen) -- (filter);
\draw[ar] (filter) -- (fwd);
\draw[ar] (fwd) -- (score);
\draw[ar] (score) -- (max);
\draw[ar,dashed,ourscol] (target.south) to[bend right=55] (score.west);

\node[right=0.3cm of max,text width=3.1cm,font=\scriptsize\itshape,text=ourscol]
  {$N$ is 2{,}048 in the released notebook, and is not stated in the paper};
\end{tikzpicture}
\caption[The cycle]{The cycle-consistency loop. A target goes in, $N$
candidate sequences come out, each is fed back through the forward task, and
the best-scoring one is reported. The dashed line is where the \Rsq{} of
Section~\ref{sec:r2trap} is computed --- \emph{within} one eight-vector.}
\label{fig:cycle}
\end{figure}

Written out:

\begin{enumerate}[leftmargin=*,itemsep=2pt]
  \item Choose an 8-vector target $\mathbf{t}$.
  \item Sample $N$ sequences from \code{GenerateSilkContent<$\mathbf{t}$>}.
  \item Discard any that do not parse as a protein.
  \item Discard any that exactly match a known silk sequence --- this is the
        paper's definition of \term{novel}.
  \item For each survivor $s$, compute $\mathbf{p} =
        \code{CalculateSilkContent}(s)$.
  \item Score $R^2(\mathbf{t}, \mathbf{p})$ across the eight entries.
  \item Report $\max$ over all survivors.
\end{enumerate}

\section{The eight property targets}

The paper picks eight targets. Three are described as random, for assessing
forward prediction (Figure 3 of the paper). Five are chosen deliberately to
span common, rare, and physically non-existent combinations (Figures 4--7).

\begin{table}[htbp]
\centering
\small
\caption[The eight targets]{The eight target vectors and the \Rsq{} the paper
reports for each. I use the labels F1--F3 and S1--S5 throughout; the paper
numbers them 1--3 and 1--5.}
\label{tab:targets}
\begin{tabular}{llcl}
\toprule
\textbf{ID} & \textbf{8-vector} & \textbf{paper \Rsq} & \textbf{rationale} \\
\midrule
F1 & [0.600, 0.204, 0.279, 0.043, 0.329, 0.127, 0.502, 0.139] & 0.8764 & random \\
F2 & [0.300, 0.200, 0.200, 0.100, 0.300, 0.100, 0.200, 0.100] & 0.8369 & random \\
F3 & [0.100, 0.200, 0.200, 0.040, 0.300, 0.100, 0.700, 0.100] & 0.6889 & random \\
\midrule
S1 & [0.250, 0.252, 0.200, 0.151, 0.310, 0.203, 0.293, 0.265] & 0.8899 & common $(E,T)$ pair \\
S2 & [0.550, 0.252, 0.750, 0.151, 0.310, 0.203, 0.293, 0.265] & 0.5640 & rare, high $E$ \\
S3 & [1.000, 0.252, 0.450, 0.151, 0.200, 0.203, 0.293, 0.265] & 0.7167 & rare, high toughness \\
S4 & [0.900, 0.252, 0.800, 0.151, 0.310, 0.203, 0.293, 0.265] & 0.7843 & \emph{non-existent} pair \\
S5 & [0.240, 0.252, 0.206, 0.151, 0.270, 0.203, 0.293, 0.265] & 0.7751 & common (strength, $T$) \\
\bottomrule
\end{tabular}
\end{table}

S4 is the interesting one scientifically: high toughness \emph{and} high
stiffness together, a combination that does not occur in the 1{,}033 measured
silks. Asking the model for it is asking for a material better than anything
nature produced. That the model returns \emph{something} is the paper's
headline capability claim.

\begin{warnbox}{An internal inconsistency worth noting}
The paper says that for each of S1--S5, the properties not being varied are set
to dataset \emph{medians}. Sets S1, S2 and S4 all carry strength $= 0.310$. Set
S3 varies toughness and $E$, so its strength should also be $0.310$ --- but both
Table~1 and the Figure~2 caption give $\mathbf{0.200}$.

Two independent places in the paper agree with each other, so this is not a
typesetting slip in one location. I kept the published value, because changing
a published input to get a better match is not reproduction. It does change
S3's target spread and therefore its \Rsq{}. Recorded as discrepancy D2.
\end{warnbox}

\section{What the paper's figures show}

Since I refer to these repeatedly, here is what each one is.

\begin{description}[leftmargin=0cm,style=nextline,itemsep=5pt]

\item[Figure 2 --- the property landscape.] A pair plot of the 1{,}033
  normalised property values: histograms on the diagonal, scatter below,
  density contours above. It establishes the distributions are roughly uniform,
  that toughness and strength are positively correlated, and it marks the five
  chosen targets in red so you can see which are common and which are out on
  the edge.

\item[Figure 3 --- forward prediction, three targets.] Bar charts. For each of
  F1--F3, eight blue bars (requested) beside eight red bars (predicted). Visual
  agreement, quantified by the \Rsq{} above each panel.

\item[Figure 4 --- the same for the five selected targets.] Same format,
  S1--S5. This is where 0.8899, 0.5640, 0.7167, 0.7843 and 0.7751 come from.

\item[Figure 5 --- do the designs look like real proteins?] For each generated
  sequence and its two closest database relatives: sequence length, molecular
  weight, instability index, isoelectric point, amino-acid composition, and
  secondary-structure fractions. The argument is that the generated sequences
  sit in the same range as natural ones.

\item[Figure 6 --- folded structures.] AlphaFold2 predictions for five
  generated sequences and one natural relative each. The paper is careful here:
  confidence scores (pLDDT) are 40--60, which is low, and it limits the claim to
  visual comparison. I did not attempt this.

\item[Figure 7 --- motifs.] Counts of specific short sequence motifs in the
  generated sequences versus their natural relatives, plus density plots of
  \emph{where} along the sequence each motif sits. The motif definitions come
  from a supplementary table which in turn derives from the Silkome paper.

\end{description}

\begin{keybox}{The one about secondary structure, which I have to flag now}
Figure 5 reports ``secondary structure fraction'', defined in its caption as
the proportion of residues that tend to be in
\begin{center}
helix (\aaseq{V I Y F W L}), \quad turn (\aaseq{N P G S}), \quad sheet (\aaseq{E M A L}).
\end{center}
Those groupings are exactly what Biopython's
\code{secondary\_structure\_fraction()} used up to version~1.81, and the caption
reproduces that function's own documentation of the time.

Biopython v1.82 changed the method, and its current docstring says why:
prior to 1.82 it \emph{``wrongly returned (Sheet, Turn, Helix) while claiming to
return (Helix, Turn, Sheet)''}. The set \aaseq{V I Y F W L} is a \emph{sheet}
propensity set; \aaseq{E M A L} is a \emph{helix} set.

So under the corrected assignment the labels in Figure 5c are exchanged. The
paper used the library as documented when it was written; the correction is
upstream and postdates publication. It matters more here than it would
elsewhere, because $\beta$-sheet content is \emph{the} central structural
quantity in spider silk (Chapter~\ref{ch:biology}), so a reader comparing
Figure 5c against the silk literature would draw the wrong conclusion about
which residues drive sheet formation. My code computes both conventions and
reports both. Recorded as discrepancy D1.
\end{keybox}

\section{What the paper claims, stated fairly}

It is worth writing down precisely what is and is not being asserted, because
much of my later analysis is about the gap between the two.

\textbf{The paper does claim:}
\begin{itemize}[leftmargin=*,itemsep=1pt]
  \item A model that generates novel spidroin sequences to order.
  \item Self-consistency between its forward and inverse directions, quantified
        by \Rsq.
  \item That generated sequences resemble natural ones on standard descriptors
        and motif content.
  \item That property combinations absent from nature can be requested and
        something is returned.
\end{itemize}

\textbf{The paper does not claim:}
\begin{itemize}[leftmargin=*,itemsep=1pt]
  \item Held-out predictive accuracy. It says plainly that fine-tuning used all
        known pairs.
  \item That the generated sequences would actually have those properties if
        synthesised. Its own limitations section says validation would require
        molecular dynamics or wet-lab work.
\end{itemize}

\begin{keybox}{The gap I set out to measure}
Both lists are accurate. But a reader encountering ``\Rsq{} $= 0.89$'' in a
property-prediction context will, by default, read it as the \emph{first} list
plus a bit of the second. The number is a maximum over an unstated $N$, of a
within-vector self-consistency measure, from a model that saw all its data.
Each of those three qualifiers is individually reasonable; together they mean
the number measures something much narrower than it appears to.

Nothing about that is misconduct or error. It is under-specification, and it is
extremely common. Measuring how much each qualifier is worth is what the four
extensions do.
\end{keybox}
